\documentclass[11pt,a4paper]{article}

\input{../../sablony/pripravy/priprava.sty}

\nastavPripravu{4. ročník}{Kombinatorika}{40}{Vlastnosti kombinačních čísel a Pascalův trojúhelník}

\begin{document}

\blok{Motivace a návaznost}

Seřaď na tabuli hodnoty $\binom n0,\ldots,\binom nn$ pro $n=0,1,2,3,4$. Nech žáky objevit zrcadlovou souměrnost a součet sousedních hodnot.

\poznamka{Předpoklady: kombinační číslo a faktoriál. Cíl: objevit symetrii, Pascalův vztah, vyplnit trojúhelník a používat jej k jednoduchým výpočtům.}

\blok{Výklad a odvození}

\begin{definiceBox}
Pro $n\in\mathbb N_0$ a $0\le k\le n$ platí \textbf{symetrie} $\binom nk=\binom n{n-k}$. Pro $n\ge2$ a $1\le k\le n-1$ platí \textbf{Pascalův vztah}
\[
\binom nk=\binom{n-1}{k-1}+\binom{n-1}{k}.
\]
\end{definiceBox}
Důkaz symetrie: volba $k$ prvků odpovídá jednoznačně volbě $n-k$ prvků, které vynecháme. Pascalův vztah: při výběru $k$ lidí rozdělíme komise podle toho, zda obsahují jednu pevně zvolenou osobu. Na krajích jsou jedničky.
\[
\begin{array}{c|l}
n&\text{řádek Pascalova trojúhelníku}\\\hline
0&1\\
1&1\quad1\\
2&1\quad2\quad1\\
3&1\quad3\quad3\quad1\\
4&1\quad4\quad6\quad4\quad1\\
5&1\quad5\quad10\quad10\quad5\quad1
\end{array}
\]
Součet hodnot v řádku $n$ je $2^n$, protože pro každý z $n$ prvků samostatně rozhodujeme, zda patří do podmnožiny.

\blok{Společně řešené příklady}

\begin{prikladReseniBox}[title={Příklad 1}]
Vypočti $\binom{10}{8}$ pomocí symetrie.
\tcblower
$\binom{10}{8}=\binom{10}{2}=45$.
\end{prikladReseniBox}



\begin{prikladReseniBox}[title={Příklad 2}]
Sestav z Pascalova trojúhelníku řádek s indexem $6$.
\tcblower
$1,6,15,20,15,6,1$.
\end{prikladReseniBox}



\begin{prikladReseniBox}[title={Příklad 3}]
Urči $\binom83$ použitím Pascalova vztahu.
\tcblower
$\binom83=\binom72+\binom73=21+35=56$.
\end{prikladReseniBox}



\blok{Samostatná práce}

\begin{zadaniReseniBox}[title={Samostatná práce 1}]
Urči $\binom91$ a $\binom99$.
\tcblower
$9$, $1$.
\end{zadaniReseniBox}



\begin{zadaniReseniBox}[title={Samostatná práce 2}]
Vypočti $\binom{11}{9}$ pomocí symetrie.
\tcblower
$\binom{11}{2}=55$.
\end{zadaniReseniBox}



\begin{zadaniReseniBox}[title={Samostatná práce 3}]
Jaký je součet sedmého řádku Pascalova trojúhelníku, pokud první řádek má index 0?
\tcblower
Jde o index $6$, součet $2^6=64$.
\end{zadaniReseniBox}



\blok{Rozšíření}

\begin{zadaniReseniBox}[title={Rozšíření 1}]
Kolik cest vede z bodu $(0,0)$ do bodu $(3,4)$, pohybujeme-li se pouze o jednotku doprava nebo nahoru?
\tcblower
Sedm kroků, z toho tři doprava a čtyři nahoru. Volíme polohy tří kroků doprava: $\binom73=35$.
\end{zadaniReseniBox}



\blok{Domácí úkol}

\poznamka{Doporučená práce doma: přibližně 15--30 minut. Úlohy jsou seřazeny od základních k náročnějším; učitel může poslední podúlohu označit jako volitelnou.}

\begin{zadaniReseniBox}[title={Domácí úkol}]

\begin{enumerate}[label=\alph*),leftmargin=*,itemsep=2pt]

\item Pomocí symetrie vypočti $\binom{12}{10}$.
\item Napiš celý řádek Pascalova trojúhelníku s indexem 5.
\item Urči $\binom72$ z Pascalova vztahu.
\item Jaký je součet všech kombinačních čísel v řádku s indexem 8?
\item Na čtvercové síti kolik monotónních cest z $(0,0)$ do $(2,3)$ vede pouze doprava a nahoru?
\item Napiš řádek Pascalova trojúhelníku s indexem 7 a ověř jeho součet.

\end{enumerate}

\tcblower

\begin{enumerate}[label=\alph*),leftmargin=*,itemsep=2pt]

\item $\binom{12}{2}=66$.
\item $1,5,10,10,5,1$.
\item $\binom61+\binom62=6+15=21$.
\item $2^8=256$.
\item $\binom52=10$.
\item $1,7,21,35,35,21,7,1$, součet $128=2^7$.

\end{enumerate}

\end{zadaniReseniBox}

\poznamka{Orientační práce v hodině: 0--15 min objevování; 15--28 min pravidla; 28--40 min aplikace; 40--45 min kontrola.}

\poznamka{Jde o obsahově nabitou hodinu; nevyžaduj složité důkazy identit. Zásadní je názornost a správný index řádku. Náročnější mřížkové úlohy ponech k rozšíření.}

\end{document}
